\section{About the Committee}

\subsection{History}

The World Trade Organization (WTO) was founded on the 1st of January, 1995, in accordance with the 1994 Marrakesh Agreement. Its precursor, the General Agreement on Tariffs and Trade (GATT), was established through a multilateral treaty between 23 countries in 1947, and it mainly delineated rules for the international trade in goods with specific focus on the gradual reduction of tariffs.\footnote{WTO. (2025a). WTO | Tariffs. Wto.org. \url{https://www.wto.org/english/tratop_e/tariffs_e/tariffs_e.htm}} The original concept of an international trade organization as a third institution to complement the World Bank and the International Monetary Fund (IMF) was excluded after negotiations came to an end in 1950. Hence, only the GATT remained as a framework administering international trade. However, owing to changing circumstances and trade policies, the rules of GATT were subject to constant amendments and appendices. The negotiations in which these modifications were discussed and decided upon are referred to as ‘rounds’.\footnote{WTO. (n.d.). Appendix 1: GATT Part IV – Trade and development [Report]. In The General Agreement on Tariffs and Trade: Part IV – Trade and Development. Commonwealth iLibrary.}

The establishment of the WTO is regarded to have marked the biggest reform of international trade since the end of World War II.\footnote{WTO (2025). WTO | the History of Multilateral Trading System. Wto.org. \url{https://www.wto.org/english/thewto_e/history_e/history_e.htm}} It superseded GATT as the most relevant organization\footnote{The GATT organization is now referred to as GATT 1947 while the current Agreement will from now on be called ‘GATT’.} and pioneered several changes to the system, such as the introduction of rules on the trade in services, for intellectual property, agriculture, textiles, and the addition of a mandatory and binding dispute settlement mechanism.\footnote{WTO. (2024). Understanding the WTO - the Uruguay round. Wto.org. \url{https://www.wto.org/english/thewto_e/whatis_e/tif_e/fact5_e.ht}} Thus, all Agreements related to the trade in goods, such as the GATT, the Agreement on Technical Barriers to Trade (TBT), the Agreement on Trade-Related Investment Measures (TRIMS), the General Agreement on Trade in Services (GATS), and the Agreement on Trade-Related Aspects of Intellectual Property Rights (TRIPS) were integrated into the Annex 1 of the Marrakesh Agreements. This brought all these treaties within the framework of the WTO, ensuring that the Ministerial Conference (MC) has the authority to amend all of them, alongside making the interpretation of all disputes related to them subject to the Dispute Settlement Understanding (DSU).\footnote{WTO. (2017). WTO | legal texts - Marrakesh Agreement. Wto.org. \url{https://www.wto.org/english/docs_e/legal_e/marag_e.htm}} As of 31 December 2024, member countries have referred 631 trade disputes for settlement to the Dispute Settlement Body (DSB).

\guidefig[WTO Members and Observers (2026). Source: BlankMap-World6]{assets/figure-01.jpeg}[0.8]

At present, 166 countries are members of the WTO, which collectively represent about 98\% of the world’s international trade. Moreover, 23 governments are qualified as observer governments, meaning that they must initiate accession negotiations within five years of being acknowledged thusly (with the Holy See being an exception). International intergovernmental organizations, such as the United Nations (UN), Food and Agriculture Organization (FAO), World Bank (WB), and World Intellectual Property Organization (WIPO) among others, can also be granted the status as an observer which enables them to follow discussions on matters directly relevant to them.\footnote{WTO. (2024a). WTO Members and Observers. Wto.org; World Trade Organization. \url{https://www.wto.org/english/thewto_e/whatis_e/tif_e/org6_e.htm}}

\subsection{Mandate}

The WTO has multiple functions: it manages a global system of trade rules, it serves as a platform for negotiating trade agreements, it resolves trade disputes between its members, and it supports the developing countries to improve their capacity to trade. All in all, the WTO is the only international organization addressing the rules of trade between nations by assisting its members in using trade as a means to raise living standards, generate employment, and enhance people’s lives.

\guidefig[Organizational Structure of the WTO. Source: WTO, 2025]{assets/figure-02.jpeg}[0.71]

At the core of the system, which is also known as the multilateral trading system, lie the WTO’s agreements, which are the legal foundations of global trade. These are negotiated and signed by a large majority of the world’s trading economies, and ratified in their parliaments. These agreements oblige governments to make their trade policies transparent by notifying the WTO about laws and measures that have been adopted by them and are currently in effect. Besides, all WTO member countries must undergo routine scrutiny of their trade policies and practices, and each of such reviews will contain reports by the WTO secretariat and the country concerned. To sum it up, WTO’s agreements facilitate a non-discriminatory trading system for its members by explicating their rights and their obligations.

To implement the rules, and therefore ensure the smooth flow of trade, the WTO’s procedure for resolving trade disputes under the DSU is indispensable. Disputes are brought to the WTO by governments if they deem their rights under the agreements to have been violated, following which judgements based on interpretations of the agreements and individual members’ commitments are given by specially appointed independent experts. Most of the time, the system urges members to resolve their conflicts through mutual consultation. However, if such a setting is unachievable, or if it proves to be unsuccessful, they can opt for a step-by-step procedure that comprises the possibility of a ruling by a panel of experts.

More than three-quarters of WTO members are developing economies or least-developed countries (LDCs). All WTO agreements have special provisions for them, which include steps to expand their trading opportunities, longer durations to fulfill commitments, and assistance in developing the infrastructure required to participate in global trade. Hundreds of training activities for developing countries are arranged by the WTO annually. It organizes technical cooperation missions to developing economies and holds many trade policy courses every year in Geneva for government officials. Prominent among the programmes by the WTO for the developing countries are the Aid for Trade initiative and the Enhanced Integration Framework (EIF), which are in effect to assist these economies in building trade capacity and using it as an engine for growth, sustainable development, and poverty reduction.

To summarize, guided by the fundamental principles of non-discrimination, predictability, transparency, sustainability, and freer trade, the WTO caters to its objective with the mandate of monitoring trade agreements, settling trade disputes, building the trade capacity of developing economies, and cooperating with other international organizations.
